\section{Introduction}
\label{sec:introduction}

A permutation of symbolic labels and a symmetry of a frozen dynamical system
are different objects.  Before arithmetic weights are assigned, tensor atoms
may be relabeled freely.  After the variables are specialized to distinct
values \(x_p=p^{-s}\), a nontrivial atom permutation changes the weighted
transfer rather than commuting with it.  Character resolution based on that
permutation therefore cannot be interpreted as a decomposition of one fixed
arithmetic operator.

Finite-group extensions supply a lawful alternative.  Their deck action lives
on a separate fiber, commutes with the base dynamics, and resolves a regular
transfer into representation blocks.  This framework is classical in
hyperbolic and symbolic dynamics
\citep{ParryPollicott1986,AdachiSunada1987,ParryPollicott1990}; the question
here is whether the signed tensor-subset grammar carries an intrinsic cocycle
whose blocks retain its exact Euler ledger.

The minimal answer is positive.  A symbol is a nonempty subset \(S\) of a
finite atom set, its scalar weight is
\((-1)^{|S|+1}x_S\), and its \(C_2\) fiber increment is
\[
\alpha(S)=|S|\pmod2.
\tag{1.1}\label{1.1}
\]
This label is subset degree, not a prime-indexed table.  Singleton symbols
switch the fiber, so the extension has nontrivial periodic data.  Deck
translations never move the variables or roofs.  The regular transfer has a
trivial block and a sign block with determinants
\[
D_+(x)=\prod_p(1-x_p),\qquad
D_-(x)=\prod_p(1+x_p).
\tag{1.2}\label{1.2}
\]
Their product
\[
\Dreg(x)=\prod_p(1-x_p^2)
\tag{1.3}\label{1.3}
\]
is the determinant of the whole two-dimensional fiber.  Equations
\eqref{1.2} and \eqref{1.3} come from one transfer and one normalization.
This same-object statement is the positive result.

It also exposes the obstruction.  The clean product follows from degree count
alone.  If a relabeling-natural one-letter cocycle is required to satisfy the
same atom-local matrix identity, coefficient comparison forces its
cardinality-\(k\) label to be the \(k\)th power of the singleton label.  The
image is cyclic.  Enlarging the finite group cannot create a clean nonabelian
escape without leaving the theorem's one-letter hypotheses.

Even the lawful \(C_2\) extension does not align its primitive cycles with
primes.  A singleton base loop closes in the full extension after two
traversals, whereas the mixed symbol \(\{p,q\}\) closes after one.  More
generally, a primitive base orbit of total subset degree \(c\) has
\(\gcd(m,c)\) primitive lifts, each traversing the base orbit
\(m/\gcd(m,c)\) times.  Atom-locality is therefore a determinant statement,
not a prime-orbit bijection.

The paper makes four claims.

\begin{enumerate}
\item \textbf{Genuine same-object factorization.}
The intrinsic parity extension is transitive at every nonempty finite cutoff,
mixing from two atoms onward, and its regular determinant factors exactly into
the two character blocks in \eqref{1.2}.
\item \textbf{Functorial one-letter rigidity.}
Under inclusion compatibility, relabeling naturality, and operator-coherent
atom locality in a faithful representation, every label is
\(a^{|S|}\).  A transitive full fiber is consequently cyclic.
\item \textbf{Primitive obstruction.}
The exact lift formula distinguishes primitive base necklaces from primitive
lifted cycles and proves that mixed lifts persist.  No orbitwise
\(p\leftrightarrow\gamma_p\) ledger survives in the whole extension.
\item \textbf{Adversarial stop.}
The factorization is a polynomial identity for arbitrary inventories.
Prime, shuffled-prime, composite, and random-rational controls all pass with
the same rate, so the identity pass-rate margin is zero.
\end{enumerate}

The exact prototype supports the algebraic theorems without fitting: all
formal determinant checks through ten atoms have zero mismatch; 300
trace-repetition rows and 350 \(C_m\) character rows are exact; the naturality
enumeration contains 72,079 tables; all 64 inventory controls and all 14 unit
tests pass.  These numbers certify implementation and boundary cases.  They do
not substitute for the proofs.

The resulting Route-A tuple is
\[
\begin{aligned}
(&\texttt{A0\_ANALYTIC\_ARITHMETIC\_ORIGIN},
\texttt{A1\_WEAK},
\texttt{A2\_ANALYTIC\_DETERMINANT},\\
&\texttt{A3\_PARTIAL\_ANALYTIC\_STRUCTURE},
\texttt{A4\_FAIL}).
\end{aligned}
\tag{1.4}\label{1.4}
\]
The A3 coordinate credits only the exact same-object Artin structure in the
honest domain; it does not borrow an external continuation theorem.
Primitive mismatch and zero control margin force
\texttt{ROUTE\_A\_REJECTED / STOP\_SCOPED / PROVES\_TOO\_MUCH}.  No Hilbert
space generator, self-adjointness argument, completed divisor, or Weil
compression is present, so Route B is locked.

\Cref{sec:classical} separates the classical Artin machinery from the scoped
claim.  \Cref{sec:source} freezes the symbolic object.
\Cref{sec:factorization} proves the exact character factors, and
\Cref{sec:rigidity} proves cyclic rigidity.  The primitive and control
obstructions occupy \cref{sec:lifts,sec:certificates};
\cref{sec:route} records the analytic and Route-A boundary.
